\section[引言]{引言（10 月 30 日）}
\srcpage{2}

本课程讲的是与 Laurent Fargues 的合作工作。

\subsection{设定}

设 $E$ 为非阿基米德局部域，即 $E=\mathbb{F}_q((t))$ 或 $\mathbb{Q}_p$ 的有限扩域。我们引入一些记号：
\begin{itemize}
  \item $\mathbb{F}_q$ 是 $E$ 的剩余域。
  \item $\pi\in\mathcal{O}_E$ 是一致化子。
  \item $G/E$ 是约化群，例如 $G=\mathrm{GL}_n$、$\mathrm{Sp}_{2n}$。（这个理论的一部分吸引力在于：它对任意约化群都完全一致地成立。）
\end{itemize}

我们关心局部投射有限群（locally profinite group）$G(E)$ 的表示论。下面简要回忆其中一些事实。

\subsection{$p$ 进群的表示}

设 $\Gamma$ 为局部投射有限群。

\begin{definition}
域 $L$ 上 $\Gamma$ 的一个光滑表示（smooth representation）\footnote{也许「局部常值」是更好的术语……}是一个 $L$-向量空间 $V$，连同映射 $\Gamma\to\mathrm{GL}(V)$，使得对一切 $v\in V$，$\mathrm{Stab}(v)\subseteq\Gamma$ 都是开子群。
\end{definition}

\begin{example}
设 $K\subset G(E)$ 为开紧子群。由于 $K$ 投射有限，它有很多有限商。设有这样的有限商 $K\twoheadrightarrow\bar K$ 以及表示 $\rho:\bar K\to\mathrm{GL}(V_0)$。于是可以作紧诱导
\[
c\text{-}\Ind_K^\Gamma(\rho)=\{f:\Gamma\to V_0\mid f\ \text{有紧支撑},\ f(gk)=f(g)k\quad\forall g\in\Gamma,\ k\in K\}.
\]
若 $L$ 的特征为零，这些对象是光滑表示范畴的紧投射生成元。
\end{example}

\begin{example}[紧诱导]
设 $K=\mathrm{GL}_n(\mathcal{O}_E)\subset\Gamma=\mathrm{GL}_n(E)$，$\bar K=\mathrm{GL}_n(\mathbb{F}_q)$，$\rho$ 为 $\mathrm{GL}_n(\mathbb{F}_q)$ 的尖点表示（cuspidal representation）。则 $c\text{-}\Ind_K^\Gamma\rho$ 是「超尖点」表示，在固定中心特征标之后不可约。
\end{example}

\begin{example}[抛物诱导]
若 $P\subset G$ 是以 $L$ 为 Levi 的抛物子群，且 $(V_0,\rho_L)$ 是 $L(E)$ 的光滑表示，则
\[
\Ind_{P(E)}^{G(E)}V_0=\{f:G(E)\to V_0\mid f(\gamma p)=f(\gamma)p\quad\forall\gamma\in G(E),\ p\in P(E)\}
\]
是 $G(E)$ 的光滑表示。超尖点表示就是不出现在这种表示中的表示。在某种模糊的意义上，$G(E)$ 的所有不可约光滑表示都由 Levi 子群的超尖点表示经抛物诱导构造出来。
\end{example}

\begin{example}[自守表示]
若 $\mathbb{E}$ 为整体域，$\mathbb{G}$ 为 $\mathbb{E}$ 上的约化群，使得 $(G,E)$ 作为 $(\mathbb{G},\mathbb{E})$ 的局部化出现，则自守形式空间
\[
\mathcal{A}\bigl(\mathbb{G}(\mathbb{E})\backslash\mathbb{G}(\mathbb{A}_{\mathbb{E}}),\mathbb{C}\bigr)
\]
是 $G(E)$ 的光滑表示。\yizhu{原文把 $(G,E)$ 与 $(\mathbb{G},\mathbb{E})$ 写成同一记号；按文意后者应为 $(\mathbb{G},\mathbb{E})$。}{A（原文 p.2）}
\end{example}

在某种意义上，对 $G(E)$ 的光滑表示范畴的研究，是研究自守形式空间的局部类比。

\srcpage{3}

\subsection{局部 Langlands 对应}

现在叙述局部 Langlands 猜想的一个多少不太精确的形式。为简单起见，设 $G$ 分裂。

\begin{conjecture}[Langlands]
考虑 $L=\mathbb{C}$ 上的表示。存在「自然的」映射
\[
\mathrm{Irrep}(G(E))/\sim\ \longrightarrow\ \mathrm{Hom}(W_E,\widehat{G}(\mathbb{C}))/\widehat{G}(\mathbb{C})\text{-conj},
\]
其中 $\widehat{G}$ 是 Langlands 对偶群，$W_E$ 是 $E$ 的 Weil 群\footnote{这是 $\mathbb{Z}\subset\widehat{\mathbb{Z}}$ 在满射 $\Gal(\bar E/E)\to\widehat{\mathbb{Z}}$ 下的原像，该满射对应 $E$ 的最大无分歧扩域。}；该映射满射，且纤维有限（纤维称为 $L$-包）。
\end{conjecture}

\begin{remark}
猜想有更精细的版本，能精确描述这些纤维。近来还有一些改进，它们猜想性地用 Artin 叠 $\mathrm{Hom}(W_E,\widehat{G}(\mathbb{C}))/\widehat{G}(\mathbb{C})$ 上的凝聚层来刻画整个范畴 $\Rep(G(E))$。
\end{remark}

以下是我最初想到的几个问题。
\begin{enumerate}[label=(\arabic*)]
  \item 这为什么说得通——$W_E$ 与 $\Rep(G(E))$ 之间如何相关？
  \item $\widehat{G}$ 从何而来？
\end{enumerate}

为进一步讨论 (2)，先回忆 $\widehat{G}$ 是什么。回忆：任意域上的分裂约化群 $G$ 由根数据 $(X,\Phi,X^*,\Phi^\vee)$ 分类。一个有趣的观察是：$(X^*,\Phi^\vee,X,\Phi)$ 也是根数据，它定义了 $\widehat{G}$。

\begin{example}[一些对偶群]
$\widehat{\mathrm{GL}}_n=\mathrm{GL}_n$，$\widehat{\mathrm{SL}}_n=\mathrm{PGL}_n$，$\widehat{\mathrm{Sp}}_{2n}=\mathrm{SO}_{2n+1}$，$\widehat{\mathrm{SO}}_{2n}=\mathrm{SO}_{2n}$。
\end{example}

\subsection{局部 Langlands 对应的例子}

\subsubsection{$G=\mathbb{G}_m$}

当 $G=\mathbb{G}_m$ 时，$G(E)=E^\times$ 交换，故 $\mathrm{Irrep}(E^\times)=\{\chi:E^\times\to\mathbb{C}^\times\}$ 是特征标的集合。对偶群为 $\widehat{G}=\mathbb{G}_m$，并且
\[
\mathrm{Hom}(W_E,\widehat{G}(\mathbb{C}))=\mathrm{Hom}(W_E,\mathbb{C}^\times)=\mathrm{Hom}(W_E^{ab},\mathbb{C}^\times).
\]
事实上，局部类域论给出同构 $E^\times\xrightarrow{\sim}W_E^{ab}$。

\subsubsection{$G=\mathrm{GL}_n$}

当 $G=\mathrm{GL}_n$ 时，对偶群为 $\widehat{G}=\mathrm{GL}_n$。

\begin{theorem}[Laumon--Rapoport--Stuhler, Harris--Taylor, Henniart]
$\mathrm{GL}_n(E)$ 的超尖点表示与 $W_E\to\mathrm{GL}_n(\mathbb{C})$ 的不可约 $n$ 维表示之间存在「自然的」双射。
\end{theorem}

这还没有给出 $\mathrm{Irrep}(\mathrm{GL}_n(E))/\sim$ 与 $\mathrm{Hom}(W_E,\widehat{G}(\mathbb{C}))$ 之间的自然双射，但用 Bernstein--Zelevinsky 可以确切说出会发生什么。

\begin{example}[自守诱导]
与 $n=1$ 的情形不同，实践中在 Galois 一侧写下对象更容易。设 $E'/E$ 为 $n$ 次扩域，$\chi':W_{E'}\to W_{E'}^{ab}\cong(E')^\times\to\mathbb{C}^\times$ 为一般特征标，则 $\varphi:=\Ind_{W_{E'}}^{W_E}(\chi')$ 是 $W_E$ 的不可约 $n$ 维表示。它应当对应 $\mathrm{GL}_n(E)$ 的一个超尖点表示 $\pi_\varphi$。直接构造 $\pi_\varphi$ 相当困难。
\end{example}

\begin{remark}
局部 Langlands 对应的许多显式例子已由 Bushnell、Kutzko、Kaletha 等人的工作给出。它极其复杂！
\end{remark}

\subsection{本课程的目标}

\srcpage{4}

\subsubsection{}

第一个目标是给出映射 $\pi\mapsto\varphi_\pi$ 的构造：从不可约表示到 $L$-参数，它对任意约化群 $G$ 都一致地成立，并且是纯粹局部的。

\begin{remark}[相关工作]
在函数域情形，这基本上已由 Genestier--Lafforgue 完成。对 $\mathrm{GL}_n$，证明走的是整体方法。此后 Jim Arthur 用扭内视（twisted endoscopy）把它推广到经典群。
\end{remark}

\subsubsection{}

第二个目标是「解释」$W_E$ 与 $\widehat{G}$ 从哪里来。

\subsubsection{}

第三个目标是把局部 Langlands 对应表述为范畴的等价这种更有结构的形式，并（本质上）构造一个方向上的函子。这使陈述精确得多：我们构造的这个函子就是等价。

\begin{remark}[系数]
这些方法对任意满足 $p\in\Lambda^\times$ 的系数环 $\Lambda$ 都适用。
\end{remark}

\subsection{粗略的想法}

我们想在「Fargues--Fontaine 曲线」上发展几何 Langlands 纲领，使用「完满空间/菱形」的几何。

我们无法详细回顾所有这些主题。参考文献包括：
\begin{itemize}
  \item \emph{Berkeley lectures on $p$-adic geometry}.
  \item \emph{Lecture notes for Montreal workshop}.
\end{itemize}

\subsection{整体图景}

我们想审视空间「$\Spec E$」。看待一个概形有几种不同的方式。
\begin{itemize}
  \item 一种是通过它的 étale 拓扑；在 $\Spec E$ 的情形，它由 $\pi_1^{\et}(\Spec E)=\Gal(\bar E/E)$ 掌控。
  \item 另一种是通过凝聚层。在这里，凝聚层就只是 $E$-向量空间。更一般地，可以考虑 $G$-主丛的群胚
  \[
  [\mathrm{pt}/G](\Spec E)=\coprod_{\alpha\in H^1(E,G)}[\mathrm{pt}/G_\alpha(E)],
  \]
  其中 $G_\alpha(E)$ 是对应于主丛 $\alpha$ 的内形式。特别地，$[\mathrm{pt}/G(E)]$ 作为开闭子叠落在其中，且 $\Rep(G(E))=\mathrm{Shv}([\mathrm{pt}/G(E)])$ 自然地嵌入 $\mathrm{Shv}([\mathrm{pt}/G](\Spec E))$。我们会发现，研究这个更大的范畴是有益的\footnote{这个想法也受到 Vogan 与 Bernstein 的强调。}。
\end{itemize}

现在我们要把这个图景稍作修改。我们想把 $\Gal(\bar E/E)$ 换成 $W_E$。对 $\mathbb{F}_q$ 上的概形 $X$，我们用带 $\Frob_q$ 作用的 $X_{\overline{\mathbb{F}}_q}$ 替换 $X$。此时 $\pi_1(X_{\overline{\mathbb{F}}_q}/\Frob_q)=\pi_1(X_{\overline{\mathbb{F}}_q})\rtimes\mathbb{Z}$ 是基本群的「Weil 形式」。

在我们的情形，这提示把
\begin{itemize}
  \item $\Spec\mathbb{F}_q((t))$ 换成「$\overline{\mathbb{F}}_q((t))/\Frob_q$」；
  \item $\Spec E$ 换成 $\Spec\breveE/\Frob_q$，其中 $\breveE$ 是 $E$ 的最大无分歧扩域的完备化。
\end{itemize}

现在改为考虑「$\Spec\breveE/\Frob_q$」的图景：
\begin{itemize}
  \item étale 侧由 $\pi_1(\Spec\breveE/\Frob_q)=W_E$ 掌控。
  \item \srcpage{5}凝聚侧，向量丛按下降就是等晶体（isocrystal）的范畴
  \[
  \Isoc_E=\{\breveE\text{-向量空间}+\text{Frob-半线性}\ \varphi:V\xrightarrow{\sim}V\}.
  \]
  这比单纯的 $E$-向量空间丰富。由 Dieudonné--Manin 定理，它分解为
  \[
  \bigoplus_{\lambda\in\mathbb{Q}}\Isoc_E^\lambda,
  \]
  其中 $\Isoc_E^\lambda$ 由斜率为 $\lambda$ 的纯等晶体组成。范畴 $\Isoc_E^\lambda$ 可以等同于 $D_\lambda$-模，其中 $D_\lambda/E$ 是不变量为 $[\lambda]\in\mathbb{Q}/\mathbb{Z}$ 的中心除代数。

  更一般地，可以考虑 $\Isoc_E$ 中 $G$-主丛的范畴。这由 Kottwitz 分类，动机来自 Shimura 簇的理论。他证明 $G$-等晶体的群胚是
  \[
  \coprod_{b\in B(E,G)}[\mathrm{pt}/G_b(E)],
  \]
  其中 $G_b$ 是 $G$ 的某个 Levi 子群的内形式。有单射 $H^1(E,G)\hookrightarrow B(E,G)$。Kottwitz 和 Kaletha 已经从纯表示论的视角意识到，同时考虑所有 $G_b$ 是有益的。
\end{itemize}

\subsection{等晶体叠}

我们想要比 $\coprod_{b\in B(E,G)}[\mathrm{pt}/G_b(E)]$ 更具几何性的东西。我们想把它提升为「$G$-等晶体叠」，对此有几种做法。

\subsubsection{}

$\Isoc_E$ 是 $E$-线性范畴，故对任意 $E$-代数 $A$ 都可以考虑 $\Isoc_E\otimes_E A$ 中的 $G$-主丛。这导出 $E$ 上的一个 Artin 叠，即
\[
\coprod_{b\in B(E,G)}[\mathrm{pt}/G_b].
\]
我们强调，这里 $G_b$ 是代数群，而这一观点对于看出 $G_b$ 可以提升为代数群至关重要。但从我们的视角看，这并不是正确的做法，因为我们追求的是 $p$ 进群的光滑表示，而不是代数群的代数表示。

\subsubsection{}

更好的想法是把 $\overline{\mathbb{F}}_q$ 换成任意（完美）$\overline{\mathbb{F}}_q$-代数 $R$\yizhu{本译本把 perfect 译作「完美」，以与 perfectoid 的定译「完满」相区分。}{I（本库 SY02-16 术语表：perfectoid=完满）}，即把
\begin{gather*}
\Spec\overline{\mathbb{F}}_q((t))/\Frob_q\ \rightsquigarrow\ \Spec R((t))/\Frob_q,\\
\Spec\breveE/\Frob_q\ \rightsquigarrow\ \Spec(W(R)\otimes_{W(\overline{\mathbb{F}}_q)}E)/\Frob_q.
\end{gather*}
我们可以在完美 $\overline{\mathbb{F}}_q$-代数上定义叠
\[
R\mapsto\bigl\{G\text{-主丛}: \Spec R((t))/\Frob_q\ \text{或}\ \Spec(W(R)\otimes_{W(\overline{\mathbb{F}}_q)}E)/\Frob_q\bigr\}.
\]
结果表明，这是一个性态良好的叠，我们称之为 $G-\Isoc$\footnote{$G-\Isoc$ 等价于 Gaitsgory、Genestier--Lafforgue、Zhu 考察过的对象「${}^{L}G/_{\sigma}{}^{L}G$」。他们给出 $D(G-\Isoc,\Ql)$ 到各分片 $D(G-\Isoc^b,\Ql)\cong D(G_b(E),\Ql)$ 的半正交分解。}。

\srcpage{6}

\begin{theorem}[Rapoport--Richartz, Caraiani--Scholze, Ivanov, Anschütz]
$G-\Isoc$ 是 $v$/arc 拓扑下的叠\footnote{这甚至强于在 fpqc 拓扑下成为叠！}。

对任意 $b\in B(E,G)$，存在局部闭子叠 $G-\Isoc^b\subset G-\Isoc$，在其上等晶体同构于 $b$，并且 $G-\Isoc^b\cong[\mathrm{pt}/G_b(E)]$\footnote{这构成一个分层，但各层不是开闭的（粘合数据很复杂）。}。
\end{theorem}

\subsubsection{}

我们仍然不会完全考虑这个叠。我们要考虑一个更解析的版本，它可以说性态更好，用完满环（perfectoid ring）代替完美代数。

为说明动机，我们指出：与 Langlands 对偶群 $\widehat{G}$ 的联系是通过 Hecke 算子。

我们把 $R((t))$ 看成一个小的去心圆盘。我们要考虑 $G$-主丛空间上的对应。

我们考虑参数化如下数据的空间：$(\Spec R((t))/\Frob_q)$ 上的两个 $G$-主丛 $\mathcal{E}_1,\mathcal{E}_2$，以及它们在某个除子 $D\subset\Spec R((t))$ 之外的同构 $\mathcal{E}_1\cong\mathcal{E}_2$。这需要截面 $\Spec R\to\Spec R((t))$，而当 $R$ 是离散环时这样的截面不存在。所以我们需要取 $R$ 为 Banach 代数。这促成了完满几何的登场。

对 $\overline{\mathbb{F}}_q$ 上任意完满 affinoid 代数 $(R,R^+)$，可以联系 Fargues--Fontaine 曲线 \mbox{$\mathbb{D}^*_{\Spa(R,R^+)}/\Frob_q$}（或混合特征下的类似对象）。于是可以考虑 Fargues--Fontaine 曲线上 $G$-主丛的模空间。

解析情形的优点是去心圆盘要几何得多。因此我们真的可以考虑 Hecke 算子
\[
\begin{tikzcd}
& \Hck_G \arrow[dl, "\mathrm{pr}_1"'] \arrow[dr, "\mathrm{pr}_2"] & \\
\Bun_G & & \Bun_G\times(\Spa\,\breveE/\phi)
\end{tikzcd}
\]
这里 $\Spa(\breveE/\phi)$ 是参数化 Fargues--Fontaine 曲线截面的空间。

$\Hck_G$ 是无穷维的，但它有有限维的层（strata）$\Hck_G^\mu$。令
\[
T_\mu:=\mathrm{pr}_{2!}\mathrm{pr}_1^*:D(\Bun_G,\Ql)\to D(\Bun_G\times\Spa\,\breveE/\phi;\Ql)=D(\Bun_G,\Ql)^{W_E}.
\]
几何 Satake 等价说，算子 $T_\mu$ 由 $\Rep(\widehat{G})$ 枚举。回忆 $D(\Bun_G,\Ql)$ 是诸 $\Rep(G_b(E);\Ql)$ 的某种汇总。结果表明，要定义附加在 $\Rep(G_b(E);\Ql)$ 上的 $L$-参数，恰恰需要这种范畴结构。
